\section{系统配置}
%% ============================================================

\subsection{config.toml 配置体系}

Codex 的配置通过 \texttt{config.toml} 文件管理，支持两个层级：

\begin{itemize}
    \item \textbf{个人默认}：\texttt{\textasciitilde/.codex/config.toml}（通过 Codex App 的 Settings $\rightarrow$ Configuration $\rightarrow$ Open config.toml 打开）
    \item \textbf{仓库特定}：\texttt{.codex/config.toml}，存放项目级别的行为配置
    \item \textbf{命令行覆盖}：仅用于一次性场景（CLI 用户）
\end{itemize}

可配置项包括：model 选择、reasoning effort、sandbox mode、approval policy、profiles、MCP setup 等。

\subsection{Sandbox 与 Approval 机制}

Codex 内置操作系统级别的沙箱机制，提供两个关键控制旋钮：

\begin{importantbox}{两个安全控制旋钮}
\begin{itemize}[nosep]
    \item \textbf{Approval mode}：决定 Codex 何时需要请求你的许可才能执行命令
    \item \textbf{Sandbox mode}：决定 Codex 能否在目录中读写，以及可以访问哪些文件
\end{itemize}
\end{importantbox}

\begin{warningbox}{新手建议：先紧后松}
如果你刚接触 coding agent，应当保持默认的严格权限。只在了解工作流、确认信任特定仓库后，才逐步放宽权限。不要一开始就给 Codex 完全的系统权限。
\end{warningbox}

\begin{knowledgebox}{CLI、IDE 和 App 共享配置}
CLI、IDE 扩展和 Codex App 共享同一套配置层级。在任一端修改配置都会影响其他端的行为。
\end{knowledgebox}

\subsection{本章小结}

通过 \texttt{config.toml} 的个人级和仓库级配置实现一致性行为。Sandbox 和 Approval 是两个核心安全旋钮，建议新手从严格权限开始，逐步放宽。

%% ============================================================
